\documentclass[10pt,a4paper]{amsart}
\usepackage[margin=19mm]{geometry}
\usepackage{amsmath,amssymb,mathtools}
\usepackage[scr=rsfs,cal=euler]{mathalfa}
\usepackage{xfrac}
\usepackage[unicode,hidelinks,bookmarks=false]{hyperref}
\newcommand{\ie}{i.e.\ }
\newcommand{\cf}{cf.\ }
\newcommand{\resp}{respectively\ }
\newcommand{\Subsection}{\subsection}
\providecommand{\nobreakdash}{\nobreak\hbox{-}}
\setlength{\emergencystretch}{2em}
\multlinegap0pt
\usepackage{amssymb}
\usepackage[shortlabels]{enumitem}
\usepackage{float}
\usepackage{tikz}
\newtheorem{proposition}{Proposition}
\renewcommand\H{{\mathcal H}}
\newcommand{\V}{\operatorname{Vol}}
\title{Distribution of lengths of closed saddle connections on moduli space of large genus translation surfaces}
\author{Shenxing Zhang}
\date{Source: arXiv:2511.12595v4 (CC BY 4.0)}
\begin{document}
\maketitle

\noindent\emph{Source and licence.} Distribution of lengths of closed saddle connections on moduli space of large genus translation surfaces. Version arXiv:2511.12595v4 (CC BY 4.0), distributed by arXiv under the Creative Commons Attribution 4.0 International licence, http://creativecommons.org/licenses/by/4.0/, as recorded in the arXiv licence metadata for this exact version and shown on the abstract page https://arxiv.org/abs/2511.12595v4. Full licence notice: \texttt{licenses/arxiv-2511.12595-CC-BY-4.0.txt}.\par\smallskip

\noindent\textit{Editorial scope.} Counted unit: the article's proposition D (source0.tex lines 597--599). The article's complete original proof of it is delivered with it (source0.tex lines 600--615, one \texttt{proof} environment of 16 lines). No mathematics is written, repaired or re-derived: the statement and the proof are the article's own lines, in the article's own notation, and no retained line is substituted or re-flowed.  No further source-local context is needed: every cross-reference of every retained line resolves inside the two delivered spans.  Nothing was omitted from any retained span, so the package carries no omission record.  No retained line contains a citation, so the wrapper carries no bibliography and no bibliography entry.  No retained line contains a figure environment, an image-inclusion command or drawing code, so no image is delivered and the package declares no asset dependency.  Editorial formatting only, in addition: an AMS \texttt{amsart} preamble that restates the article's own declarations and macros used by the retained lines, namely the environments \texttt{proposition} on separate counters, together with the standard packages those lines need.  The retained environments print in this document as Proposition 1 in order of appearance, whereas the article prints the same items with the numbering of its own full document.  The complete native source of the article as distributed in the arXiv e-print for this exact version is bundled unchanged as \texttt{sources/189209/source0.tex}.
\begin{proposition}\label{D}
$$\V(\mathcal{D})=O(\frac{1}{g})\V(\H_g(2^{g-1}))$$
\end{proposition}

\begin{proof}
For some $(X,f,\omega)\in \mathcal{D}$, there is a closed saddle connection $\gamma$ goes through at least $g^{2/3}$ handles, and $\ell_\omega(\gamma)\leq \frac{b}{\sqrt{g}}$.
Consider the geodesic representation $\overline{f(H_i)}$ of each handle $f(H_i)$ and denote $\gamma_i=\overline{f(H_i)}\cap \gamma$.

Since $\ell_\omega(\gamma)\leq \frac{b}{\sqrt{g}}$, there exists some $i$ such that $\ell_\omega(\gamma_i)\leq\frac{b}{g^{7/6}}$.
Now $\overline{f(H_i)}$ is a translation surface with geodesic boundary, and $\gamma_i$ is a nontrivial curve in $H_1(\overline{f(H_i)},\partial{\overline{f(H_i)}})$.
Consider the geodesic representation of $\gamma_i$ in $\overline{f(H_i)}$, it consists of some saddle connections, so it implies there is at least a saddle connetions
of length less than $\frac{b}{g^{7/6}}$.
From Siegel-Veech formula, the volume of $\mathcal{D}$ can be estimated by
$$\V(\mathcal{D})\leq (g-1)(g-2)c \frac{b^2}{g^{7/3}} \V(\H_g(2^{g-1}))=O(\frac{1}{g})\V(\H_g(2^{g-1})),$$
where the factor $(g-1)(g-2)$ is the ways to choose two zeros from $g-1$ zeros and 
$c$ is the Siegel-Veech constant connecting two double zeros on $\H_g(2^{g-1})$, which
is $O(1)$.


\end{proof}

\end{document}
